% =========================================================
% 1.4 JUSTIFICACIÓN
% El \section{} de esta sección está en el chapter.tex del capítulo;
% sus referencias van en seccion1_4.bib (misma carpeta).
% =========================================================

En el diseño de rutas peatonales seguras, los sistemas actuales suelen fundamentarse en modelos reactivos impulsados por estadísticas de delitos pasados. Sin embargo, en México este enfoque presenta un sesgo metodológico crítico al omitir la cifra negra, la cual supera el 90\% de los actos delictivos no reportados. Al priorizar únicamente las áreas con carpetas de investigación, las rutas generadas ignoran el riesgo latente en zonas no documentadas, generando la necesidad de recurrir a variables auxiliares para superar esta limitación. La Criminología Ambiental y la Teoría de Patrones Delictivos sugieren que el crimen no es aleatorio, sino que está estructurado por un 'telón de fondo ambiental' compuesto por la infraestructura física y las actividades rutinarias \cite{brantingham1995criminality}.  Al integrar estas variables espaciales, es posible transitar de un modelo reactivo a uno predictivo, capaz de calcular rutas seguras evaluando la vulnerabilidad del entorno, mitigando así la dependencia exclusiva de los índices de denuncia.

El presente Trabajo Terminal propone solucionar este vacío, mediante el desarrollo de un prototipo de aplicación móvil que implemente un modelado de riesgo socio-espacial. A diferencia de los mapas de calor tradicionales, este sistema ponderará factores de diseño ambiental estáticos (como callejones o baldíos) y dinámicos (alumbrado público, horarios comerciales), así como incidencias de crímenes para generar rutas que minimicen la exposición al riesgo a nivel de calle. Mientras que las soluciones actuales nos mencionan estadísticas como “Aquí asaltaron ayer”, nuestro sistema advierte “Este entorno propicia el delito”. De esta manera, no solo se provee al usuario de “la mejor ruta”, sino que también se le propician datos específicos sobre el estado de las calles por las que transita. Asimismo, no se limita a buscar una ruta segura, pues es muy probable que las personas no estén dispuestas a desviarse por mucho de la ruta más corta. Así, con la culminación del trabajo presentado, se tendrá una combinación de un modelo preventivo para generar rutas seguras, acompañado de un análisis de distancia para asegurar una ruta “óptima” o “balanceada”.

Frente a la necesidad de transitar por zonas desconocidas en la alcaldía Gustavo A. Madero, la aplicación busca devolver a los transeúntes la tranquilidad para poder recorrer su ciudad con libertad mediante la mitigación del sentimiento de vulnerabilidad. Según datos de la Encuesta Nacional de Seguridad Pública Urbana (ENSU), tan solo en el primer trimestre de 2025, el 40.5\% de los entrevistados reconoció haber cambiado hábitos en cuanto a caminar de noche por los alrededores de su vivienda y el 25.5\% cambió rutinas relacionadas con visitar parientes o amistades \cite{inegi2025ensu}. Al proveer información sobre las condiciones de seguridad ambiental y la incidencia delictiva de las rutas disponibles, el proyecto empodera al usuario y lo llama a tomar decisiones informadas, brindando la posibilidad de reducir su riesgo de victimización, una ruta a la vez.

Finalmente, la materialización de este proyecto exige la aplicación de conocimientos adquiridos en el mapa curricular, destacando Análisis y Diseño de Algoritmos para la lógica de enrutamiento, Inteligencia Artificial para el modelado de riesgo socio-espacial, Bases de Datos para la gestión topológica de la información, y Desarrollo de Aplicaciones Móviles Nativas para la interfaz de usuario. En términos de factibilidad técnica y económica, el proyecto es altamente factible debido a la arquitectura “pre-procesamiento” donde la utilización de modelos de IA y la API de Google Maps (Street View y Places) se destinarán de manera exclusiva al poblado inicial de la base de datos. Esta estrategia elimina la dependencia y los costos asociados a llamadas de API en tiempo real durante la operación de la aplicación. Con el polígono de estudio delimitado a la alcaldía Gustavo A. Madero, se asegura que la recolección de datos y el modelado puedan ejecutarse dentro de los tiempos establecidos y a su vez, permite que periódicamente se realicen actualizaciones sobre el índice de riesgo de las calles.
